\section{From arbitrary matrices to the unique equality class}\label{sec:proof}
The left kernel of a matrix in our rank-constrained class supplies the
comparison functional. A positive stress gives a trace bound whose
equality condition determines the matrix entries.

\begin{lemma}[Left-kernel trace bound]\label{lem:trace}
Let $A\in\mathbb C^{N\times N}$ have diagonal one, rank at most $N-2$,
and coherence $\mu=\max_{i\ne j}|A_{ij}|$. Choose
$Y\in\mathbb C^{2\times N}$ with rows in the left kernel of $A$ and
normalization $YY^*=NI_2/2$, so that $YA=0$.
For every $K\succ0$ with $\tr K=2$,
\begin{equation}\label{eq:trace-bound}
 N(1+1/\mu)\le\Gamma_K(Y).
\end{equation}
If equality holds and all off-diagonal entries of $H=Y^*KY$ are nonzero,
then $A_{ij}=-\mu H_{ij}/|H_{ij}|$ for every $i\ne j$.
\end{lemma}
\begin{proof}
Take two orthonormal columns $Q$ in $\ker A^*$ and put
$Y=\sqrt{N/2}\,Q^*$. Then $HA=0$, $H\succeq0$, and $\tr H=N$.
Also $\mu>0$, since the identity has rank $N$. Expanding
$\tr(HA)=0$, separating the diagonal terms, and taking real parts gives
\[
 N=-\sum_{i\ne j}\Rea(H_{ij}A_{ji})
 \le\mu\sum_{i\ne j}|H_{ij}|=\mu(\Gamma_K(Y)-N).
\]
If equality holds, every ordered term with $H_{ij}\ne0$ attains both
its modulus and real-part bounds. Thus
$A_{ji}=-\mu\overline{H_{ij}}/|H_{ij}|$, which is the stated formula
after exchanging the indices.
\end{proof}

\subsection{The quantitative comparison used in the proof}
For $Y$ as in Lemma~\ref{lem:trace}, write $y_i=r_i u_i$ and let $x_i\in S^2$
be the Bloch vector of $u_i$. If $r_i=0$, choose $u_i$ arbitrarily at
this stage. Put
\[
 a=\frac{2\sqrt3-1}{8},\quad k=\frac{6\sqrt3-9}{4},\quad L=\sum r_i,
 \qquad g(t)=(1-\sqrt t)(3/4+k\sqrt t)(\sqrt t-1/\sqrt3)^2.
\]
The radial, moment and angular deficits are
\begin{align}\label{eq:deficits}
 u&=a(N^2-L^2),&
 v&=\frac k4\left\|\sum_i r_i x_ix_i^T-\frac L3I\right\|_F^2,&
 w&=\sum_{i,j}r_ir_jg((1+x_i\cdot x_j)/2).
\end{align}
They are nonnegative and satisfy
$u+v+w\le\alpha N^2-\Gamma_I(Y)$, by the exact polynomial identity
in Section~\ref{app:global}. Applying Lemma~\ref{lem:trace} with $K=I$
to a competitor with $\mu\le U_N=(N\lambda_{\rm bal}-1)^{-1}$ gives
\begin{equation}\label{eq:budget}
 u+v+w\le M,
 \qquad M=\begin{cases}1/2,&N\text{ odd},\\11/19,&N\equiv2\pmod4.
 \end{cases}
\end{equation}
The scalar candidate estimates establishing these budgets are included
in Section~\ref{app:candidate}.

\begin{proposition}[Localization of every column]\label{prop:localization}
Suppose $N\ge180$, $4\nmid N$, $YY^*=NI_2/2$, and
$\Gamma_I(Y)\ge N^2\lambda_{\rm bal}$. There is a partition of all
$N$ columns into four nonempty groups with positive counts $m_a$.
Set $\delta_a=m_a-N/4$ and $C=\sum\delta_a^2$. Up to permutation and
simultaneous sign reversal, its possible profiles are
\begin{center}
\begin{tabular}{ccc}
\toprule
Residue & $C$ & $(\delta_1,\delta_2,\delta_3,\delta_4)$\\
\midrule
$N$ odd & $3/4$ & $(-1,-1,-1,3)/4$\\
$N$ odd & $11/4$ & $(-5,-1,3,3)/4$\\
$N\equiv2\pmod4$ & $1$ & $(-1,-1,1,1)/2$\\
$N\equiv2\pmod4$ & $3$ & $(-3,1,1,1)/2$\\
\bottomrule
\end{tabular}
\end{center}
On this same partition let $\bar x_a$ be the group means and
$S=\sum_a\sum_{i\in a}\|x_i-\bar x_a\|^2$.
Then $S\le31E/N$, where $E=\sum_{i,j}g((1+x_i\cdot x_j)/2)$.
The simultaneous total and pointwise estimates
\eqref{app:global:LU-LP}--\eqref{app:global:physical-inputs}
hold in one common preliminary alignment to the repeated tetrahedral
frame $F=(f_{a(i)})_i$, where $a(i)$ is the group containing column $i$
and the $f_a$ are unit tetrahedral spinors.
\end{proposition}
\begin{proof}
Section~\ref{app:global} keeps the shared budget $u+v+w\le M$
through anchor selection and partitions every column. Its centered count
estimate gives $C<5$ in residue two and $C<19/4$ in the odd residues,
which imply exactly the listed profiles. The forward refinement on
this same partition gives the scatter and simultaneous polar estimates.
The finite inequalities used in these steps are specified in
Section~\ref{app:certificates}.
\end{proof}

The scales explain the role of the local step. The shared deficit is
$O(1)$, uniformly in $N$, and the count variance $C$ is bounded.
Localization and the candidate estimates place both $Y$ and the
candidate $Z$ with the same profile within $O(N^{-1/2})$ of $F$ in the
Frobenius norm. The triangle inequality and minimization over the allowed
alignments therefore give $d=\|Y-Z\|_F=O(N^{-1/2})$ in the closest gauge.
The separate transport of pointwise bounds to that gauge is proved in
Section~\ref{app:local:angular}. The relative normal component is at most
$d/\sqrt{2N}=O(N^{-1})$. The horizontal reference
curvature has a negative gap bounded away from zero; candidate
perturbations and the higher-order loss tend to zero. This explains why
the local comparison becomes strict for large $N$. To obtain the
explicit threshold $180$, we bound both terms on the localized regions
and certify the resulting inequalities by rational arithmetic.

\begin{proposition}[Conditional local comparison]
\label{prop:conditional-local}
Fix one of the localized profiles, its candidate $Z$ and its positive
stress $K$. Put a competitor $Y$ with $YY^*=NI_2/2$ in the
closest unitary and individual-phase gauge to $Z$, allowing conjugation
of the candidate, and set $d=\|Y-Z\|_F$. Suppose all validity conditions
in Section~\ref{app:local}, including the denominator, positivity, phase,
overlap and closest-gauge transport conditions, hold throughout a closed
region containing this competitor. Let $c_m,c_r$ be valid lower bounds
for the curvature gaps in the mean and residual directions, and $J,L$
valid upper bounds for
the mixed and pure normal terms, measured in the original Frobenius
metric as in \eqref{eq:main-quadratic}. Let
$A_{\rm rem},B_{\rm rem}\ge0$ supply one bound for the full Taylor
remainder using the common tangent fraction defined below. Define
$g_{\rm normal}$ as the negative of the maximum of the quadratic upper
bound in \eqref{eq:main-quadratic} over the normalized components, with
the normal fraction restricted by the bound for this region. If
\[
 g_{\rm normal}-\frac{A_{\rm rem}+
          \sqrt{A_{\rm rem}^2+B_{\rm rem}^2}}2\ge\varepsilon>0,
\]
then
\begin{equation}\label{eq:conditional-comparison}
 \Gamma_K(Y)\le N^2\lambda_p-\varepsilon N\|Y-Z\|_F^2.
\end{equation}
\end{proposition}
\begin{proof}
The case $d=0$ is immediate. For $d>0$, define the Lagrangian
$\mathcal L(Y)=\Gamma_K(Y)-N\lambda_p\tr(KYY^*)$ and its half-Hessian
$Q(H)=\tfrac12D^2\mathcal L(Z)[H,H]$, with $K$ fixed and derivatives
taken over the underlying real vector space. Put $\Delta=Y-Z$ and
$R_n=(\Delta\Delta^*/N)Z$. The horizontal component
$\Xi=\Delta+R_n$ splits orthogonally into cluster means $H_{\rm m}$ and
zero-sum residuals $H_{\rm r}$. Set $h=\|H_{\rm m}\|_F/d$,
$r=\|H_{\rm r}\|_F/d$, and $n=\|R_n\|_F/d$. Under the stated validity
conditions, the half-Hessian estimate is
\begin{equation}\label{eq:main-quadratic}
 \frac{Q(\Delta)}{Nd^2}\le-c_mh^2-c_rr^2+2Jnh+Ln^2,
 \quad h^2+r^2+n^2=1,\quad n\le d/\sqrt{2N}.
\end{equation}
The exact maximum of this quadratic upper bound over the allowed normal
fraction, including both branches, is $-g_{\rm normal}$; see
\eqref{app:local:normal-max}.

The remainder uses the distinct physical decomposition
$y_i=(1+t_i)z_i+\xi_i$, where $t_i\in\mathbb R$ and $z_i^*\xi_i=0$.
For the common tangent fraction
$\theta=(\sum_i\|\xi_i\|^2)^{1/2}/d\in[0,1]$, its complete upper bound is
$Nd^2[A_{\rm rem}\theta^2+B_{\rm rem}\theta\sqrt{1-\theta^2}]$.
Lemma~\ref{app:local:loss} bounds this by
$Nd^2(A_{\rm rem}+\sqrt{A_{\rm rem}^2+B_{\rm rem}^2})/2$.
The first variation of $\mathcal L$ vanishes at $Z$, and the constraint
term $N\lambda_p\tr(KYY^*)$ has the same value at $Y$ and $Z$ because
their covariances agree. Adding the quadratic and remainder bounds proves
\eqref{eq:conditional-comparison}. Section~\ref{app:local} derives the
bounds and their validity conditions; the next proposition certifies
those conditions on the regions supplied by localization.
\end{proof}

\begin{proposition}[Finite verification and propagation]
\label{prop:finite-acceptance}
The exact verification statement in Section~\ref{app:certificates:acceptance}
covers all closed budget cells of the four profiles. The local endpoints
are $180$ for $C\in\{3/4,11/4,3\}$ and $182$ for $C=1$.
Each cell is either excluded by a strict counting inequality or lies in a
closed region where all validity conditions of
Proposition~\ref{prop:conditional-local} hold and the quadratic gap exceeds
the remainder bound by more than $1/10000$.
For each fixed signed deviation vector $(\delta_a)$, these inequalities
remain valid for every larger integer $N$ for which
$m_a=N/4+\delta_a$ are integers.
\end{proposition}
\begin{proof}
The rational records and the conditions checked by the verifier are given in
Section~\ref{app:certificates:acceptance}.
Lemma~\ref{app:local:propagation} extends the verified inequalities by
fixing valid endpoint coefficients and bounding the reciprocal parameter
over the full subsequent interval.
\end{proof}

\begin{corollary}[Localization and strict comparison]
\label{prop:comparison}
Under the hypotheses of Proposition~\ref{prop:localization}, use its
partition and the closest gauge to the candidate with the same profile.
Then
\begin{equation}\label{eq:strict-comparison}
 \Gamma_K(Y)\le N^2\lambda_p-
                  \frac{N}{10000}\|Y-Z\|_F^2.
\end{equation}
\end{corollary}
\begin{proof}
Combine Propositions~\ref{prop:localization},
\ref{prop:conditional-local} and \ref{prop:finite-acceptance}.
Odd integers in the range start at $181$; the first actual residue-two
integer is $182$. Thus the two different local endpoints leave no
integer gap for nondivisible $N\ge180$.
\end{proof}

\subsection{Completion of the proof and equality cases}
We now prove Theorem~\ref{thm:main}. For $4\nmid N$, take any complex
$N\times N$ matrix $A$ with diagonal entries one, rank at most $N-2$,
and coherence $\mu\le U_N$. Choose $Y$ as in Lemma~\ref{lem:trace}. The closest
gauge is legitimate for a general matrix: replacing $Y$ by $UYD$ is
accompanied by $A\mapsto D^*AD$. It preserves $YA=0$, covariance,
diagonal, rank and coherence. Select the stress $K$ for this profile \emph{after}
this alignment. Lemma~\ref{lem:trace}, Corollary~\ref{prop:comparison}
and strict balancing give
\[
\begin{aligned}
 N(1+1/\mu)&\le\Gamma_K(Y)\\
 &\le N^2\lambda_p-\frac N{10000}\|Y-Z\|_F^2\\
 &\le N^2\lambda_{\rm bal}\le N(1+1/\mu).
\end{aligned}
\]
Every inequality is an equality. Therefore $Y=Z$ in the chosen gauge,
the counts are balanced, and $\mu=U_N$. All stress overlaps are nonzero,
so the ordered equality condition in Lemma~\ref{lem:trace} forces
\[
 A=(1+\mu)I-\mu S_p=G_p.
\]
Undoing the allowed transformations gives exactly the equality class
of the theorem.

For $4\mid N$, the classical scalar bound is sharp. For completeness,
the deficit identity at zero forces $r_i=1$, vanishing moment defect,
and $g((1+x_i\cdot x_j)/2)=0$ for every pair. Thus distinct Bloch
directions have inner product $-1/3$. At most four such directions exist.
Since $r_i=1$, tightness gives $\sum x_i=0$; together with the zero
moment defect this gives $\sum l_il_i^T=NI_4/4$, where
$l_i=(1,\sqrt3 x_i)^T/2$. The lifted directions
therefore span $\mathbb R^4$, so there are exactly four. They form an
orthonormal basis of $\mathbb R^4$, and tightness forces equal counts
$N/4$. Hence the kernel consists of equally repeated tetrahedral spinor
lines. Equality in the ordered trace bound with $K=I$ again fixes
$A=G_p$. This also proves the full matrix equality statement in the
divisible case.

Proposition~\ref{prop:construction} supplies an attaining rank-exactly
$N-2$ Gram matrix. Since the rank-at-most problem is a relaxation of
both other problems, all three optimal values coincide. Finally, two
full-rank vector factors with the same Gram matrix differ by a common
unitary on $\mathbb C^{N-2}$, completing the vector classification.
